%  LaTeX support: latex@mdpi.com
%=================================================================
%  Built against the official MDPI template.
%
%  OVERLEAF: upload the MDPI template package first so that
%  Definitions/mdpi.cls and its logo files are present, then add this
%  file plus the three figures:
%      fig_validation.pdf
%      fig1_ceiling_definition.pdf
%      fig_inversion_wide.pdf
%
%  SUGGESTED SPECIAL ISSUE:
%      "Thermal and Optical Remote Sensing: Evaluating Urban Green
%       Spaces and Urban Heat Islands in a Changing Climate"
%       (section: Urban Remote Sensing)
%       https://www.mdpi.com/journal/remotesensing/special_issues/RS_Urban_Green_UHI
%  VERIFY THE DEADLINE YOURSELF -- mdpi.com blocks automated fetching,
%  so the scope match is checked but the closing date is not.
%
%  BEFORE SUBMITTING:
%    * Complete the GenAI disclosure in \acknowledgments so it matches
%      what you actually did. MDPI requires this.
%    * MDPI requires FULL author lists; entries below marked "et al."
%      must be expanded from the original sources.
%    * Reference [heatmat2025] has a placeholder author field. Fill it
%      in, or delete the entry and its row in Table 1.
%=================================================================

\documentclass[remotesensing,article,submit,oneauthor]{Definitions/mdpi}

%=================================================================
% MDPI internal commands - do not modify
\firstpage{1}
\makeatletter
\setcounter{page}{\@firstpage}
\makeatother
\pubvolume{1}
\issuenum{1}
\articlenumber{0}
\pubyear{2026}
\copyrightyear{2026}
\datereceived{ }
\daterevised{ }
\dateaccepted{ }
\datepublished{ }

%=================================================================
\Title{Separating Model Failure from Data Limits in Urban Heat Island
Simulation: A Land-Cover Data-Ceiling Diagnostic Applied Across Five
Continents}

%\newcommand{\orcidauthorA}{0000-0000-0000-000X} % add \orcidA{} after the name

\Author{Abhiram Donepudi}

\AuthorNames{Abhiram Donepudi}

\address{%
$^{1}$ \quad Westwood High School, Round Rock, TX 78681, USA;
abhiram.donepudi@gmail.com}

\corres{Correspondence: abhiram.donepudi@gmail.com}

\abstract{When an urban heat island (UHI) model transfers poorly to a new
city, two explanations are normally indistinguishable: the model may fail,
or the input land cover may not separate thermal regimes well enough for
any land-cover-driven model to succeed there. We introduce a land-cover
data-ceiling control that separates them. It is deliberately independent
of any particular simulator, requiring only a city boundary, a land-cover
raster and one satellite scene---no model run, no calibration and no
meteorology. We define the ceiling as the multiple correlation between
MODIS land surface temperature and the land-cover channels a model
consumes, and show the simpler univariate alternative (impervious fraction
alone) is unsound: it goes negative in arid cities even where the
multivariate ceiling exceeds $R = +0.45$. We further show that a ceiling
bounds a model only when its \emph{receptive field} matches too. A ceiling
regressed on each cell's own land cover does not bound a model that reads
across space: our own simulator exceeds Austin's per-cell ceiling on its
held-out days, and matching the ceiling's spatial reach restores the
bound. Across 49 US cities a built--vegetated thermal-gap heuristic
predicts the ceiling (Spearman $\rho = +0.69$, $p = 3 \times 10^{-8}$).
Extending to 28 cities on five continents under a uniform city definition
shows the arid inversion is climatic rather than an artifact of an all-US
sample---but also that it is a pixel-scale effect: all eleven inverting
cities recover to ceilings between $+0.47$ and $+0.91$ once neighbourhoods
are visible, no city in the sample remains data-limited, and the
arid/non-arid ceiling difference ceases to be detectable ($p = 0.018$ at
pixel scale, $p = 0.115$ matched). What limits arid cities is the scale at
which land cover is read, not the city.}

\keyword{urban heat island; land surface temperature; MODIS; class
separability; model transferability; surface energy balance; open data;
thermal remote sensing}

%=================================================================
\begin{document}

\section{Introduction}

A recurring problem in UHI model evaluation is that two distinct failure
modes are conflated when a model is transferred to a new city. The model
may genuinely fail to capture that city's thermal structure, or the input
land-cover data may not separate thermal regimes well enough for
\emph{any} land-cover-driven model to succeed. The two demand opposite
responses---the first calls for a better model, the second for better
inputs or an admission that the city is not tractable with these
channels---yet they produce the same symptom, a low correlation against
observed land surface temperature (LST), and the literature offers no
standard way to tell them apart.

This paper introduces a \emph{land-cover data-ceiling control} that makes
the distinction explicit. The ceiling is a per-city bound on the skill
recoverable from a given set of land-cover inputs, computed from
observations alone. Because it involves no simulator, no calibration and
no meteorological forcing, it is not a property of the model under test.

A ceiling is only a bound relative to a stated model class, and we show
that the class must be matched in two respects rather than one. The
channel count is the obvious requirement. The less obvious and more
consequential one is the \emph{receptive field}: a ceiling regressed on
each cell's own land cover does not bound a model that reads across
space, and our own simulator exceeds such a ceiling in its reference
city. Matching the ceiling's spatial reach to the model's restores the
bound, and changes which cities are judged data-limited.

We exercise the diagnostic on \textsc{austin-twin}, a physics-based slab
energy-balance simulator driven entirely by zero-authentication open data
\cite{oke1982energetic,voogt2003thermal,masson2000physically}. The
simulator is a vehicle here rather than the contribution. Existing UHI
tools span two extremes. High-fidelity mesoscale models---WRF-SLUCM, TEB,
the Urban Weather Generator \cite{bueno2013urban,kusaka2001simple}---
resolve canopy processes but require proprietary or authenticated inputs
and hours of compute per scenario. Machine-learning approaches
\cite{lyu2022integrated,hoang2025geospatial,bhamjee2024detection}
validate well against historical observations but cannot predict the
effect of an intervention outside their training distribution, which is
precisely what urban planning asks of them. \textsc{austin-twin} occupies
the middle: a 2D grid model running in 10.8~ms per 48-hour simulation on
a desktop CPU, with an explicit counterfactual engine and fully open
inputs. Table~\ref{tab:novelty} positions it; the data-ceiling column has
no prior entry.

A second, methodological motivation concerns reproducibility. Most urban
climate pipelines depend on authenticated inputs---reanalysis behind a
registration wall, licensed imagery, institutional data agreements. That
dependency is invisible in a published result and fatal to independent
verification: a reader who cannot obtain the inputs cannot check the
outputs. Every input used here is retrievable by anonymous HTTP request,
which is a constraint on the design rather than a convenience.

The contribution is not a new physical formulation but four reusable
methodological results, each independent of this simulator:

\begin{enumerate}
\item A \emph{land-cover data-ceiling control}, together with the two
specification requirements that make it a bound rather than a number.
It must be defined across all channels a model consumes---the univariate
alternative goes negative in arid cities---and its receptive field must
match the model's, since a per-cell ceiling does not bound a model that
reads across space. We pair it with a built--vegetated thermal-gap
heuristic validated across 49 cities, and an NDVI channel screen shown to
be a predictable, general rule rather than a single-city result.
\item A demonstration that receptive field is not a technicality: our own
simulator exceeds Austin's per-cell ceiling on its held-out days, and
matching the ceiling's spatial reach both restores the bound and reverses
which of two transfer-failure cities is judged data-limited.
\item Evidence that the arid inversion of the land-cover--temperature
relationship is \emph{climatic} rather than continental, established on 28
cities across five continents under a city definition uniform enough to
support the comparison---together with the finding that the effect is
definition-sensitive in the American Southwest, which we report rather
than omit.
\item A \emph{decay-constrained two-stage calibration} quantifying the
trade-off between MODIS pattern fidelity and intervention spatial
realism.
\item An \emph{evapotranspiration ablation} showing that a standard linear
evapotranspiration (ET) term contributes nothing to spatial-pattern skill
at 500~m, regardless of formulation.
\end{enumerate}

% TODO: reference heatmat2025 below still has a placeholder author field.
\begin{table}[H]
\caption{Positioning relative to existing UHI simulation and prediction
systems. ``Zero-auth'' means all inputs are retrievable by anonymous HTTP
request; ``counterfactual'' means the system evaluates an intervention
outside the training distribution via an explicit causal mechanism;
``ceiling diagnostic'' means a data-ceiling control separating model
failure from input class-separability limits.\label{tab:novelty}}
\begin{tabularx}{\textwidth}{lCCCCC}
\toprule
\textbf{System} & \textbf{Scale} & \textbf{Zero-Auth} &
\textbf{Counterfactual} & \textbf{Sat.\ Calib.} & \textbf{Ceiling Diag.} \\
\midrule
WRF-SLUCM/TEB \cite{kusaka2001simple,masson2000physically} & Meso/canyon & No & Yes & Rare & No \\
Urban Weather Generator \cite{bueno2013urban} & Neighborhood & No & Yes & No & No \\
ML (NGBoost/CNN/GeoFM) \cite{lyu2022integrated,hoang2025geospatial} & City--regional & Partial & No & Implicit & No \\
HeatMat \cite{heatmat2025} & District & Partial & No & Yes & No \\
\textbf{This work} & City & Yes & Yes & Yes & Yes \\
\bottomrule
\end{tabularx}
\end{table}

%=================================================================
\section{Materials and Methods}

\subsection{Study Areas and City Definitions}

Austin, Texas is the reference city, with Phoenix (AZ), Miami (FL) and
Denver (CO) as transfer targets chosen to span the humid--arid range. The
wide sweep covers 49 US cities (55 attempted, 6 dropped for insufficient
cloud-free coverage). The global extension covers 28 cities on five
continents.

City delimitation requires care once the sample leaves the United States.
OpenStreetMap administrative units are not internationally comparable:
querying ``Riyadh, Saudi Arabia'' returns a 375,478~km$^2$ province,
``Perth'' a 6~km$^2$ council area, and ``Alice Springs'' 1~km$^2$. All
three return a valid geometry and would silently produce a number. For
the global extension, therefore, every city---US and non-US alike---is
defined as a 25~km-radius disc about its centre (1960~km$^2$, within the
US sample's own range of administrative areas). Including US cities under
this identical definition provides a comparability bridge, so that any
US/non-US difference cannot be an artifact of how a city was delimited.
Each city is sampled in its own warm season, with southern-hemisphere
cities drawing on an austral-summer date pool.

\subsection{Data}

All inputs are retrieved anonymously, without API keys, accounts or
licences. Austin's administrative boundary comes from OpenStreetMap.
Land cover uses ESA WorldCover v200 at 10~m \cite{zanaga2022esa},
aggregated to a 500~m UTM grid ($\sim$2896 cells for Austin) as class
fractions rather than a majority class, giving three channels: impervious,
vegetation and water fraction. MODIS Aqua MYD11A1 v6.1 daily LST is
retrieved through the Microsoft Planetary Computer STAC API
\cite{microsoft_pc}. The Aqua afternoon overpass ($\sim$13:30 local) is
used deliberately, as it falls near peak surface temperature where the UHI
signal is strongest. LST is read in its native scaled-integer form with
the fill value masked before conversion, and reprojected by
nearest-neighbour resampling so that fill values cannot bleed into valid
neighbours. ERA5 hourly reanalysis \cite{hersbach2020era5} is read from
Google's analysis-ready cloud-optimised (ARCO) Zarr mirror, which is
anonymously readable; this replaces the registration-gated Climate Data
Store API. MODIS NDVI (MOD13A1, 500~m) supports the channel screen.

\subsection{Surface Energy Balance Simulator}

Each grid cell evolves under a slab energy balance,
\begin{equation}
\frac{dT_i}{dt} = Q_{\mathrm{solar}} - Q_{\mathrm{rad}}
                - Q_{\mathrm{ET}} + Q_{\mathrm{diff}},
\label{eq:slab}
\end{equation}
with per-class solar absorption, Newtonian radiative cooling, linear
evapotranspiration $Q_{\mathrm{ET}} = k_{\mathrm{ET}} f_{\mathrm{veg}}
\max(T_i - T_{\mathrm{air}}, 0)$, and a Fickian diffusion term. Water
cells use an explicit mixed-layer heat capacity and a fixed latent-heat
flux rather than a single relaxation parameter. The timestep satisfies
the Courant--Friedrichs--Lewy condition for every diffusivity in the
search range.

\subsection{Calibration}

Five coefficients are tuned by differential evolution
\cite{storn1997differential} against three training MODIS days, with two
held out for testing. The \emph{ship configuration} uses a two-stage
procedure: fix $D = 1.0$~m$^2$~s$^{-1}$, matching the literature
tree-cooling half-distance \cite{ziter2019scale}, then refit the
remaining four coefficients under ARCO-ERA5 forcing. Ship coefficients are
transferred to Phoenix, Miami and Denver under three schemes: zero-shot;
Scheme~B, re-optimising $k_{\mathrm{ET}}$ only; and Scheme~A,
re-optimising all four coefficients. All configurations being compared use
identical optimiser settings and seed, so that differences between them
are not confounded with differences in search effort.

\subsection{The Land-Cover Data Ceiling}

A \emph{data-ceiling control} establishes the maximum correlation a model
can achieve in a given city given the information its inputs carry. A
ceiling is only meaningful relative to a stated model class, and we find
that two properties of that class must be matched before the bound holds:
how many channels the model reads, and how far it reads across space.

\subsubsection{Channel Count}

The first is straightforward. A model consuming three channels can exceed
a one-channel correlation by construction, so a ``percent of ceiling
achieved'' computed against the univariate $r(\mathrm{LST},
\text{impervious})$ is not interpretable. Worse, in arid cities the
univariate quantity goes negative even where the three-channel fit exceeds
$R = +0.45$, because desert vegetation runs hotter than built surfaces
there; a negative denominator makes the ratio meaningless. We therefore
define the ceiling over all channels the model consumes, as a multiple
correlation $R$.

\subsubsection{Receptive Field}
\label{sec:receptive}

The second is less obvious and, we find, more consequential. A multiple
correlation regresses LST on each cell's \emph{own} land cover. A model
with any spatial reach does not work that way: our simulator diffuses
heat, so a modelled cell depends on its neighbours' land cover as well.
Convolutional and foundation-model approaches
\cite{lyu2022integrated,hoang2025geospatial} have receptive fields wider
still. Such a model has access to information a per-cell ceiling never
offered its estimator, and can exceed that ceiling with nothing wrong.

This is not hypothetical. Our own shipped configuration exceeds Austin's
per-cell ceiling on the very days it is scored on
(Section~\ref{sec:rfresult}), which is what led us to the distinction.

We therefore also report a \emph{receptive-field-matched} ceiling,
computed by giving the estimator mask-aware neighbourhood averages of the
same three channels at $\sigma = 1, 2, 4$ cells (roughly 0.5--2~km),
bracketing the shipped diffusion length of $\approx$590~m. Smoothing is
applied to land cover, which is a known input everywhere in the city, and
never to LST, so no label information is introduced.

\subsubsection{Estimation}

Both ceilings are estimated two ways: as a multiple correlation, and as
the out-of-fold correlation of a gradient-boosted regressor
\cite{friedman2001greedy}, which removes the linearity assumption. The
latter uses \emph{spatial} block cross-validation over contiguous
$6 \times 6$-cell tiles. This is not optional. Adjacent 500~m cells are
strongly autocorrelated, so random folds train on a test cell's own
neighbours and the score inflates toward unity, manufacturing the
conclusion that land cover explains everything.

Because single-day ceilings are noisy, we report multi-day means
throughout (5--12 cloud-free dates per focus city; median 6 across the
49-city sample).

\subsection{Statistical Methods}

Model skill is reported as the spatial correlation $r$ and RMSE of
\emph{mean-subtracted} fields. The slab formulation carries no
absolute-temperature anchor, so only the spatial pattern is meaningfully
scored; an RMSE computed on raw fields additionally carries a constant
offset the model does not attempt to predict, and would understate
agreement for reasons unrelated to pattern skill. We state this
explicitly because recomputing RMSE the obvious way yields a substantially
different number.

Because neighbouring grid cells are strongly autocorrelated, per-city
confidence intervals come from a spatial block bootstrap resampling
contiguous $6 \times 6$-cell tiles ($B = 2000$); a cell-wise bootstrap
would understate intervals by treating each cell as independent.
Comparisons between two configurations evaluated on the same cities use a
\emph{paired} bootstrap, so the interval describes the difference rather
than being inferred from two overlapping marginal intervals.

City-level associations---gap against ceiling, aridity against
ceiling---use Spearman rank correlation, since neither relationship is
assumed linear and single extreme cities should not dominate. Group
comparisons (arid versus humid; arid US versus arid non-US) use the
Mann--Whitney $U$ test, reported with group sizes, because these samples
are small enough that a non-significant result may reflect limited power
rather than absence of an effect.

%=================================================================
\section{Results}
\label{sec:results}

\subsection{Calibration and Uncertainty}

Table~\ref{tab:bootstrap_ci} reports spatial block-bootstrap 95\%
confidence intervals. The ship configuration achieves $r = +0.754$
[$+0.718$, $+0.785$], trading 0.10 in $r$ against an unconstrained
single-stage fit ($r = +0.850$) in exchange for physically realistic
diffusion and intervention decay lengths (Section~\ref{sec:disc}). ERA5
forcing contributes $+0.007$ [$+0.004$, $+0.009$] over synthetic
forcing---small but distinguishable from zero, with the benefit being
mainly methodological, in that it permits absolute-temperature
comparability without anomaly normalisation.

\begin{table}[H]
\caption{Block-bootstrap 95\% confidence intervals, held-out test $r$ and
RMSE. Both statistics are computed on mean-subtracted
fields.\label{tab:bootstrap_ci}}
\begin{tabularx}{\textwidth}{lCC}
\toprule
\textbf{Configuration} & \textbf{Test} {\boldmath$r$} \textbf{[95\% CI]} & \textbf{RMSE (}{\boldmath$^{\circ}$}\textbf{C)} \\
\midrule
Single-stage, synthetic forcing & $+0.850$ [$+0.823$, $+0.872$] & 1.09 \\
Two-stage, synthetic forcing & $+0.747$ [$+0.712$, $+0.778$] & 1.48 \\
\textbf{Ship (ARCO-ERA5)} & {\boldmath$+0.754$} [$+0.718$, $+0.785$] & \textbf{1.40} \\
\bottomrule
\end{tabularx}
\end{table}

Figure~\ref{fig:validation} shows the simulated field against the MODIS
observation for a held-out day, in both absolute and anomaly form. The
model reproduces the hot urban core and the cool riparian corridor, but
runs a mean 3.66~$^{\circ}$C cold in absolute terms (panel~c); the
anomaly panels (d--f) show what the reported statistics actually score.
The residual is not spatially uniform, a point we return to in
Section~\ref{sec:disc}.

\begin{figure}[H]
\includegraphics[width=\textwidth]{fig_validation.pdf}
\caption{Simulated versus observed LST for Austin, 21 August 2024, a
held-out day. (\textbf{a}) Simulated peak-heat field under the ship
configuration; (\textbf{b}) MODIS Aqua observation; (\textbf{c}) raw
difference, which carries the $-3.66$~$^{\circ}$C mean offset;
(\textbf{d},\textbf{e}) the same two fields after mean subtraction;
(\textbf{f}) residual of the anomalies, which is what the reported
$r = +0.739$ and anomaly RMSE of 1.49~$^{\circ}$C are computed from
($n = 2868$ cells). Panels (a,b) share one colour scale, as do
(d,e).\label{fig:validation}}
\end{figure}
\unskip

\subsection{The Ceiling Must Match the Model's Receptive Field}
\label{sec:rfresult}

A per-cell ceiling does not bound a spatial model, and the clearest
demonstration is our own. On the two held-out days of
Table~\ref{tab:bootstrap_ci}, the shipped configuration scores
$r = +0.754$ while Austin's per-cell multivariate ceiling on those same
days is $+0.693$. The model stands at 109\% of its own ceiling.

Relaxing linearity does not account for this. A gradient-boosted
regressor on the same three per-cell channels, under spatial block
cross-validation, reaches $+0.681$ on those days---marginally
\emph{worse} than the linear fit, and still exceeded by the simulator.
Across the four focus cities the per-cell nonlinear gain is small and
inconsistently signed (Austin $-0.024$, Denver $-0.014$, Miami $+0.047$,
Phoenix $+0.110$). Whatever the simulator is exploiting, it is not
nonlinearity.

It is spatial reach. Giving the estimator the neighbourhood-averaged
channels of Section~\ref{sec:receptive} raises Austin's ceiling on those
days to $+0.883$, placing the model at 85\% of it. The violation
disappears. Table~\ref{tab:receptive} gives the multi-day picture.

\begin{table}[H]
\caption{Data ceilings under three specifications, multi-day means. The
context column adds mask-aware neighbourhood averages of the same three
channels at $\sigma = 1, 2, 4$ cells. Gain is the context ceiling minus
the per-cell nonlinear one.\label{tab:receptive}}
\begin{tabularx}{\textwidth}{lCCCC}
\toprule
\textbf{City} & \textbf{Linear, per-Cell} & \textbf{Nonlinear, per-Cell} &
\textbf{Nonlinear, Context} & \textbf{Gain} \\
\midrule
Austin & $+0.613$ & $+0.589$ & $+0.789$ & $+0.200$ \\
Phoenix & $+0.271$ & $+0.381$ & $+0.689$ & $+0.307$ \\
Denver & $+0.437$ & $+0.423$ & $+0.756$ & $+0.333$ \\
Miami & $+0.841$ & $+0.887$ & $+0.952$ & $+0.065$ \\
\bottomrule
\end{tabularx}
\end{table}

The gain is not an artifact of smoothing crossing fold boundaries. At
block sides of 6, 12 and 24 cells---the last four times the smoothing
reach---Austin gains $+0.202$, $+0.208$ and $+0.170$, and Phoenix
$+0.310$, $+0.252$ and $+0.290$.

Two consequences follow. First, a percent-of-ceiling figure is only
interpretable once the ceiling's receptive field matches the model's;
quoted against a per-cell ceiling, it can exceed 100\% for any model with
spatial reach. Second, and more practically, the gain is largest in the
arid cities and smallest in Miami, which we return to in
Section~\ref{sec:disc}.

\subsection{Multi-City Transfer and the Ceiling Diagnostic}

Table~\ref{tab:transfer} presents the transfer matrix. Against per-cell
ceilings Miami reaches 99\% under Scheme~A ($r = +0.832$ against
$+0.841 \pm 0.025$), while Phoenix and Denver reach 28\% and 29\%
($+0.271 \pm 0.085$ and $+0.437 \pm 0.085$).

Read against receptive-field-matched ceilings, however, the picture
changes materially. Phoenix and Denver have ceilings of $+0.689$ and
$+0.756$ (Table~\ref{tab:receptive}), so the achieved fractions become
11\% and 17\%. Both cities retain substantial unrecovered headroom, and
the earlier reading---that Denver's higher ceiling leaves room for
calibration to exploit while Phoenix's does not---is not supported.
Re-optimisation still degrades Phoenix ($+0.077 \rightarrow +0.033$) and
roughly doubles Denver ($+0.056 \rightarrow +0.126$), but an absence of
recoverable signal in Phoenix is no longer the explanation: the signal is
present, and this simulator is not capturing it.

We report this rather than the more convenient earlier framing because it
is what the matched ceilings show, and because it sharpens what the
diagnostic is for. The ceiling bounds what is recoverable; it does not
diagnose \emph{why} a particular model fails to recover it.

\begin{table}[H]
\caption{Multi-city transfer, held-out test $r$ over 2 days. The ceiling
is the multi-day mean $R$(LST; impervious, vegetation,
water).\label{tab:transfer}}
\begin{tabularx}{\textwidth}{lCCC}
\toprule
\textbf{Scheme} & \textbf{Denver} & \textbf{Phoenix} & \textbf{Miami} \\
\midrule
Austin ship (reference) & $+0.754$ & $+0.754$ & $+0.754$ \\
Zero-shot & $+0.056$ & $+0.077$ & $+0.724$ \\
Scheme B, 3 days & $+0.063$ & $+0.071$ & $+0.776$ \\
Scheme A, 3 days & $+0.126$ & $+0.033$ & $+0.832$ \\
\midrule
\textbf{Ceiling} & {\boldmath$+0.437 \pm 0.085$} & {\boldmath$+0.271 \pm 0.085$} & {\boldmath$+0.841 \pm 0.025$} \\
\bottomrule
\end{tabularx}
\end{table}

Figure~\ref{fig:ceiling} shows why the univariate definition fails. Every
city lies on or above the $y = x$ line, because the multiple correlation
is greater than or equal to the univariate one by construction. A
three-channel model can therefore exceed a one-channel ``ceiling'' with
nothing wrong. Five cities fall left of $x = 0$, where a percentage of
ceiling is undefined.

\begin{figure}[H]
\includegraphics[width=\textwidth]{fig1_ceiling_definition.pdf}
\caption{Univariate against multivariate data ceiling for the 49-city US
sample. Points on or above the diagonal show that the multiple
correlation is at least the univariate one by construction, so a
multi-channel model can exceed a univariate ``ceiling'' without any
defect. Cities left of $x = 0$ have a negative univariate ceiling, for
which a percent-of-ceiling figure is undefined.\label{fig:ceiling}}
\end{figure}
\unskip

\subsection{The Thermal-Gap Heuristic Across 49 Cities}

The built--vegetated LST gap---mean LST over impervious cells minus mean
LST over vegetated cells---is the physical driver behind the ceiling, and
is far cheaper to compute. It validates across the 49-city sample: against
the sound multivariate ceiling, Spearman $\rho = +0.69$
($p = 3 \times 10^{-8}$).

A univariate proxy target gives a more flattering $\rho = +0.84$,
confirmed significantly higher by a paired bootstrap
($\Delta\rho = +0.15$, [$+0.03$, $+0.29$]). We report both because the
difference is diagnostic: part of the univariate figure's apparent
strength comes from a shared arid sign-flip affecting both variables, not
from gap predictiveness as such.

The heuristic generalises across arid-exclusion subsets ($p < 10^{-4}$)
with one exception, the humid half ($n = 24$, $\rho = +0.36$,
$p = 0.08$). We leave this open: at this sample size, underpowering and a
genuine concentration of the association in arid cities are not
separable, and we prefer to say so rather than assert the more convenient
reading.

The association is not an artifact of measuring the ceiling at pixel
scale. Recomputing all 49 ceilings at matched receptive field gives
$\rho = +0.716$ ($p = 7 \times 10^{-9}$) against $+0.694$ published; a
paired bootstrap resampling cities together puts the difference at
$+0.022$ with 95\% CI $[-0.085, +0.142]$. The heuristic therefore
predicts recoverable signal rather than merely per-pixel recoverable
signal, which is what makes it usable as a screen for models of any
receptive field. Under matched ceilings no US city is data-limited
either---the lowest is Spokane at $+0.335$, the median $+0.785$, and none
is negative---and all five cities that invert under OpenStreetMap limits
recover (Boise $-0.331 \rightarrow +0.747$, Albuquerque
$-0.308 \rightarrow +0.716$, Tucson $-0.254 \rightarrow +0.602$, Phoenix
$-0.095 \rightarrow +0.662$, Billings $-0.034 \rightarrow +0.418$).

In arid climates the land-cover--temperature relationship \emph{inverts}
rather than merely weakening. Natural desert vegetation runs hotter than
irrigated, shaded built surfaces, driving the univariate ceiling negative
(Figure~\ref{fig:inversion}).

\begin{figure}[H]
\includegraphics[width=\textwidth]{fig_inversion_wide.pdf}
\caption{The inversion rendered as surface relief, where height and
colour both encode the plotted variable. The built core rises as a
plateau in \emph{both} cities' impervious channel (\textbf{a},\textbf{c};
range 0--1), but the temperature surface rises with it in Austin
(\textbf{b}; 38--48~$^{\circ}$C) and sinks into a basin in Riyadh
(\textbf{d}; 53--57~$^{\circ}$C), where the surrounding desert is hotter
than the city. Nothing in the land-cover input distinguishes the two
cases, which is precisely why a univariate land-cover ceiling reports
that no model can work in the arid one. Single representative cloud-free
dates; Riyadh under the 25~km disc definition.\label{fig:inversion}}
\end{figure}
\unskip

\subsection{Is the Inversion Climatic or Continental?}

The 49-city sample is entirely American, so ``arid'' is confounded with
``American Southwest'' and the inversion could reflect a continental
artifact---a difference in construction practice, irrigation policy or
land-use zoning rather than climate. Adding a sixth Southwest city cannot
resolve this. We therefore repeat the diagnostic on 28 cities across five
continents, under the uniform disc definition described in
Section~2.1.

The inversion is climatic. Eleven cities show negative univariate
ceilings and nine of them lie outside the US (Table~\ref{tab:arid}),
spanning Arabia, the Levant, North Africa, the Atacama, Iberia, Mexico
and central Australia. Arid and
humid/temperate ceilings differ sharply (median $-0.032$ against
$+0.615$; Mann--Whitney $p = 2.4 \times 10^{-4}$), while arid US and arid
non-US cities do \emph{not} differ detectably ($p = 0.075$)---the
comparison an all-US sample cannot make. The gap heuristic replicates on
this independent sample ($\rho = +0.92$ univariate, $+0.57$
multivariate), as does the aridity link (mean LST against ceiling,
$\rho = -0.69$).

\begin{table}[H]
\caption{Every city with a negative multi-day mean univariate ceiling
under the uniform 25~km disc definition, on five continents. The context
column is the receptive-field-matched ceiling of
Section~\ref{sec:receptive}: no inverting city remains
low.\label{tab:arid}}
\begin{tabularx}{\textwidth}{llCCC}
\toprule
\textbf{City} & \textbf{Region} & \textbf{Univariate} &
\textbf{Context} & \textbf{Gap (}{\boldmath$^{\circ}$}\textbf{C)} \\
\midrule
Amman & Levant & $-0.429$ & $+0.858$ & $-2.52$ \\
Riyadh & Arabia & $-0.342$ & $+0.612$ & $-1.56$ \\
Marrakesh & North Africa & $-0.156$ & $+0.743$ & $-1.18$ \\
Antofagasta & South America & $-0.124$ & $+0.687$ & $-3.38$ \\
Hermosillo & North America & $-0.116$ & $+0.686$ & $-0.99$ \\
Madrid & Europe & $-0.102$ & $+0.468$ & $-0.51$ \\
Billings \textsuperscript{1} & North America & $-0.101$ & $+0.640$ & $-0.61$ \\
Alice Springs & Australia & $-0.096$ & $+0.559$ & $-3.11$ \\
Cairo & North Africa & $-0.082$ & $+0.912$ & $+1.08$ \\
Yuma \textsuperscript{1} & North America & $-0.067$ & $+0.866$ & $-0.48$ \\
Mendoza & South America & $-0.032$ & $+0.859$ & $-0.57$ \\
\bottomrule
\end{tabularx}

\noindent{\footnotesize{\textsuperscript{1} United States; the remaining
nine lie outside it.}}
\end{table}

\subsubsection{A Robustness Caveat We Report Rather Than Omit}

Inversion in the American Southwest is sensitive to how the city is
delimited. Of the five US cities that invert under OpenStreetMap city
limits, four (Phoenix, Tucson, Albuquerque, Boise) do \emph{not} invert
under the disc definition; only Billings does. The hyper-arid non-US
cities invert decisively under either definition (Antofagasta gap
$-3.38$~$^{\circ}$C, Alice Springs $-3.11$~$^{\circ}$C), so the mechanism
itself is real, but the US examples sit near the sign boundary and should
not be read as definition-independent.

Palm Springs is the converse control. It shows no inversion ($+0.255$)
and the sample's largest positive gap ($+6.34$~$^{\circ}$C), in the same
desert as Yuma, because its landscape is irrigated---the mechanism
behaving exactly as predicted in a case selected to test it.

\subsubsection{The Inversion Is a Pixel-Scale Effect}
\label{sec:pixelscale}

Applying the receptive-field-matched ceiling of
Section~\ref{sec:receptive} to all 28 cities of the global extension
changes what the inversion means, though not whether it is real.

Every city gains, from $+0.019$ (Chicago) to $+0.593$ (Boise), and the
size of the gain is predicted almost entirely by how low the per-cell
ceiling was (Spearman $\rho = -0.869$, $p = 2 \times 10^{-9}$) rather
than by aridity directly ($\rho = +0.310$, $p = 0.11$). Arid cities gain
more (median $+0.328$ against $+0.167$, Mann--Whitney $p = 0.005$)
because they are the cities whose per-cell ceilings were low, not because
aridity acts on the gain through some separate route.

Every inverting city recovers. All eleven cities in
Table~\ref{tab:arid} reach a context ceiling between $+0.468$ and
$+0.912$. More strikingly, no city anywhere in the 28-city sample is
data-limited once the ceiling can see neighbourhoods: the lowest context
ceiling is Windhoek's $+0.463$ and the median is $+0.779$. The spread
across cities collapses from $0.858$ at pixel scale to $0.526$ with
context, so much of the apparent variation in data readiness is an
artifact of the scale at which readiness was measured.

Consequently the arid/non-arid ceiling difference, which is clear at
pixel scale (median $+0.372$ against $+0.632$, $p = 0.018$), is no longer
detectable once receptive fields are matched ($+0.727$ against $+0.828$,
$p = 0.115$). The thermal-gap heuristic survives this, weakened but
significant: $\rho = +0.571$ against the linear ceiling, $+0.540$
per-cell and $+0.483$ against the context ceiling ($p = 0.009$).

The inversion is therefore a property of pixel-scale land cover rather
than a limit on modelling arid cities. The built core still separates
thermally from the desert around it; what fails is the assumption that
individual built and vegetated cells separate.

\subsection{Channel Screening Without a Simulator}

Because the ceiling depends only on input data, candidate channels can be
screened before any model is built. Adding MODIS NDVI raises the ceiling
predictably across the 49-city sample: the gain correlates with ceiling
height ($\rho = -0.55$, $p = 4.9 \times 10^{-5}$) and with aridity
($\rho = +0.50$, $p = 2.7 \times 10^{-4}$), with 27 of 49 cities robustly
positive and a median gain of $+0.06$.

Phoenix---the irrigated-desert city that motivated adding the channel in
the first place---ranks 4th of 49 ($+0.22$), behind Bakersfield, Spokane
and Boise. The channel is therefore a general, aridity-driven screening
rule rather than a single favourable city's result, which is the
distinction that determines whether such a screen is worth applying
elsewhere.

\subsection{Evapotranspiration Ablation}

Replacing the linear ET term with canonical Penman--Monteith reduces test
$r$ from $+0.755$ to $+0.416$. Leaf-area-index saturation and
vapour-pressure-deficit decoupling---the two mechanisms we isolated---
together explain only 21\% of this regression. Stomatal resistance $r_s$
recovers 61\% at its optimum ($r_s = 3200$~s~m$^{-1}$, roughly
$32\times$ a reference grass value, and not physically defensible), but
even there Penman--Monteith ($r = +0.686$) underperforms simply
\emph{removing} the ET term ($r = +0.758$; Table~\ref{tab:pm}).

No formulation we tested---including one tuned to a physically
indefensible parameter specifically to maximise its own score---recovers
what deleting the term already achieves. This reframes the question from
``which ET formulation is best'' to ``does this linear ET term do useful
work at 500~m at all''. On this evidence it does not.

\begin{table}[H]
\caption{Evapotranspiration formulation ablation. Each configuration is
independently recalibrated, so the comparison is not confounded by
coefficients fitted for a different model.\label{tab:pm}}
\begin{tabularx}{\textwidth}{lC}
\toprule
\textbf{Configuration} & \textbf{Test} {\boldmath$r$} \\
\midrule
No ET term & {\boldmath$+0.758$} \\
Linear ET (reference) & $+0.755$ \\
PM, $r_s = 3200$~s~m$^{-1}$ (optimum, indefensible) & $+0.686$ \\
PM, canonical ($K = 3$, $r_s = 100$~s~m$^{-1}$) & $+0.416$ \\
\bottomrule
\end{tabularx}
\end{table}

\subsection{Counterfactual Scenarios}

Under ship-configuration ERA5 forcing, a $+20\%$ canopy increase produces
a mean $\Delta T = -0.74$~$^{\circ}$C, a citywide slope of
$-0.44$~$^{\circ}$C per 10\% canopy, within the literature band
\cite{zhou2025tree,wri2026cooling}. A river greenway yields
$-0.25$~$^{\circ}$C and cool roofs $-0.07$~$^{\circ}$C
\cite{santamouris2014cool,wang2022berlin}; suburban densification, run as
a control in the opposite direction, yields $+0.18$~$^{\circ}$C.

These means are averaged over every (cell, timestep) pair in the 48-hour
run. The convention matters: the same scenario gives
$-1.52$~$^{\circ}$C at the final frame alone, and per-frame values swing
between roughly $-0.3$ and $-1.5$~$^{\circ}$C across the diurnal cycle,
so a cooling figure quoted without its averaging window is ambiguous by a
factor of two. The spatial pattern is likewise non-uniform: cooling
concentrates where there was impervious surface to convert, and water
cells warm slightly relative to baseline.

Two quantities fall outside their literature bands. Local per-patch
cooling ($-0.59$~$^{\circ}$C per 10\%) is 1.5--5$\times$ below field
measurements \cite{krayenhoff2021albedo}, and the edge-referenced
half-cooling distance ($\approx$330~m) sits just outside the 100--300~m
band. Both reflect the diffusion trade-off analysed below rather than
independent defects.

%=================================================================
\section{Discussion}
\label{sec:disc}

\subsection{What Became of the Data Limit}
\label{sec:dichotomy}

This paper began from a dichotomy: a low transfer score means either that
the model failed or that the input data could not support any model. The
diagnostic was built to tell those apart. It is worth stating plainly what
happened when we applied it, because the result is not the one the framing
anticipates.

Across 49 US cities under administrative boundaries and 28 cities under a
uniform disc definition, \emph{no city is data-limited} once the ceiling
is specified to match the model's receptive field. The lowest matched
ceiling anywhere is $+0.335$; the median is near $+0.78$ in both samples.
Every city-definition whose univariate ceiling is negative---sixteen of
them, eleven in the global sample and five in the US, the cases that
motivated the arid-inversion result---recovers to a strongly positive
matched ceiling. The second branch of the dichotomy, as measured by land
cover, is close to unoccupied.

We read this as a positive finding rather than a null one. Low
land-cover-versus-temperature correlations are real and are routinely
interpreted as evidence that a city is intrinsically hard to model. Our
results say that interpretation is usually a specification artifact: the
correlation was measured at one scale and the model operates at another.
The practical question shifts from \emph{whether} land cover can support a
model in a given city to \emph{at what scale} it does, which is a question
with an actionable answer everywhere we looked.

\subsection{The Real Data Limit Is Observability, Not Separability}

The dichotomy is not empty, however. It is occupied by a failure mode the
ceiling does not measure.

Of 35 cities attempted in the global extension, seven yielded no usable
warm-season observation at all---not a weak ceiling, but zero cloud-free
days clearing our coverage threshold. The US sweep lost a further six of
55 the same way. These cities have no ceiling because they have no data,
and the diagnostic is silent on them by construction, since it is computed
from scenes that exist.

The losses are not confined to one climate. Bangkok, Manila and Lagos are
humid tropical, but Dubai, Kuwait City and Jodhpur are arid, where dust
and haze trigger the same cloud masking as cloud does. Roughly one city in
seven across both samples was lost this way, and the loss falls hardest on
exactly the places where heat exposure is most severe and where modelling
capacity is scarcest.

This also conditions our own headline. ``No city is data-limited'' is
measured among cities we could observe, and observability is not random
with respect to climate. A complete data-readiness assessment therefore
needs two axes, not one: whether land cover separates thermal regimes at
some accessible scale, and whether the thermal observations needed to
check that exist in the first place. We report the first here and can only
flag the second, which we regard as the more urgent gap for any attempt to
extend this diagnostic globally.

\subsection{The Diffusion--Fidelity Trade-Off}

No single Fickian diffusivity $D$ satisfies both MODIS spatial
autocorrelation and literature cooling-decay scales. The best $r$
achievable inside the 100--300~m band is $+0.67$ at $D = 0.20$, a cost of
$-0.19$ in $r$ against the unconstrained optimum ($r = +0.85$ at
$D \approx 12.6$). This is roughly twice the cost a naive, uncalibrated
estimate would suggest, and it means a model tuned purely for pattern
correlation will produce intervention footprints too diffuse to use for
planning.

\subsection{Advection Is Not a Drop-In Fix}

Replacing diffusion with ERA5 wind-driven advection destroys skill
($r = -0.11$ at 25\% strength). The reason is structural rather than a
matter of tuning: a spatially uniform wind translates the entire surface
pattern downwind, whereas the physical process transports heat in the air
layer above the surface. A coupled air layer, not a one-line
substitution, is required \cite{mitraka2025prototypes}.

\subsection{Class Separability as an A Priori Screen}

Because the ceiling depends only on input data, it can be computed for a
candidate city before any modelling effort is committed. Given
Section~\ref{sec:dichotomy}, what that buys is not a shortlist of cities
worth modelling---on present evidence that is most of them---but the scale
at which each city's land cover becomes informative. For an operational
deployment across many cities this still reorders the work: establish the
informative scale first, then calibrate at that scale rather than
discovering the mismatch through a poor transfer score.

What the screen does \emph{not} do is diagnose why a given model fails.
Our own results make the distinction concrete. Under per-cell ceilings
Phoenix looks data-limited and Denver looks recoverable; under
receptive-field-matched ceilings both have ample headroom
(Section~\ref{sec:rfresult}), and Phoenix's poor transfer becomes a
statement about this simulator rather than about Phoenix. A low ceiling
is evidence that no model of the matched class will do well. A high
ceiling with a low score localises the problem to the model, which is
useful, but it does not say which part.

\subsection{Scale, Not Aridity, Is What Limits These Cities}

Section~\ref{sec:pixelscale} changes the practical reading of the arid
result. Arid cities are not intrinsically resistant to land-cover-driven
modelling; they are poorly served by \emph{pixel-scale} land-cover models
specifically. The remedy indicated is a wider receptive field rather than
additional channels, and the two are not interchangeable: widening the
receptive field raises Phoenix's ceiling by $+0.31$, whereas adding an
NDVI channel raises it by $+0.22$ at pixel scale and leaves the
pixel-scale limitation in place.

It is worth being precise about what aridity does here, because the
obvious reading is wrong. Every city in both samples gains from a wider
receptive field---all 49 US cities, minimum $+0.026$---so this is not a
correction that applies to arid cities alone. What varies is the size of
the gain, and the gain tracks how low the per-cell ceiling was (Spearman
$\rho = -0.707$ in the US sample, $-0.869$ globally) far more closely than
it tracks temperature ($+0.208$, $p = 0.15$; $+0.310$, $p = 0.11$).
Splitting the US sample at median LST, hot and cool halves gain
indistinguishably ($+0.212$ against $+0.225$, $p = 0.73$). Arid cities
gain more in the global sample because they are the cities whose per-cell
ceilings are low, not because aridity acts through some separate route.

This reframes what a readiness assessment should report. A single
per-city number presumes the ceiling is a property of the city. It is
better understood as a property of the city \emph{and} the scale at which
its land cover is read, and the useful output for a practitioner is the
scale at which land cover becomes informative there---which is wider than
one pixel essentially everywhere, and widest where a pixel-scale
assessment would have declared the city hopeless.

\subsection{Structure in the Residual Bounds What Calibration Can Buy}

The residual in Figure~\ref{fig:validation}f is spatially structured
rather than noise-like. The model reproduces the hot core but places it
slightly west of where MODIS observes it, leaving a coherent dipole. This
is a positional error, not random misfit, and it implies that
$r = +0.75$ is limited by something further coefficient tuning cannot
remove. Three candidate causes, in increasing cost to test, are the
diffusion trade-off smearing the core, an absorption composite too coarse
to place it precisely, and a peak-frame selection offset from the Aqua
overpass time. Establishing which dominates would bound what calibration
can still achieve, and we flag it as the most informative next diagnostic.

\subsection{Limitations}

A $\sim$7~$^{\circ}$C absolute bias, arising from the simplified longwave
scheme, couples only weakly to skill ($\Delta r = 0.02$ across a
3~$^{\circ}$C bias sweep), but it does mean absolute temperatures should
not be read as calibrated predictions. The 500~m grid smooths
highway-corridor hot streaks that finer imagery resolves. Per-city
ceilings should be read as multi-day estimates, since single-day values
can be actively unrepresentative.

Two results remain open at the present sample sizes: the humid-subset
association of Section~3.4, and the definition-sensitivity of the US arid
cities.

A third limitation is one of specification. Our matched ceiling uses
three Gaussian scales on three channels, chosen to bracket the shipped
diffusion length, rather than a search over receptive fields. The matched
ceiling is therefore at least as high as we report and may be higher,
which makes the direction of every conclusion here safe and the
magnitudes lower bounds. A model class with a wider or differently shaped
receptive field---a deep convolutional stack, for instance---would need
its own matched ceiling before a percent-of-ceiling figure against it
would mean anything.

%=================================================================
\section{Conclusions}

A land-cover data ceiling separates two failure modes that a transfer
score alone conflates, but only when it is specified to match the model it
is meant to bound. Two requirements matter. The ceiling must span the
channels a model consumes---the univariate alternative goes negative in
arid cities, where a percent-of-ceiling figure becomes undefined---and its
receptive field must match the model's, since a ceiling regressed on each
cell's own land cover does not bound a model that reads across space. We
found the second the hard way: our own simulator clears Austin's per-cell
ceiling on the days it is scored on.

Five results transfer beyond the simulator used to demonstrate them.
First, the ceiling diagnostic itself, validated across 49 cities and made
actionable through a cheap thermal-gap heuristic and a predictable NDVI
channel screen. Second, evidence that the arid inversion is climatic
rather than continental, established on 28 cities across five continents
under a uniform city definition---while being explicit that the American
Southwest cases are definition-sensitive.

Third, and the result we did not expect, that the inversion is a
pixel-scale effect. Across both samples every inverting city recovers
once the ceiling is allowed to read neighbourhoods---to between $+0.47$
and $+0.91$ globally, and $+0.42$ to $+0.75$ in the US---and no city in
either sample remains data-limited (global minimum $+0.463$, US minimum
$+0.335$). The arid/non-arid ceiling difference ceases to be detectable
globally ($p = 0.018$ at pixel scale, $p = 0.115$ matched). What limits
these cities is the scale at which land cover is read, not the city, and
the correction applies everywhere rather than to arid cities alone: all
49 US cities gain, with the size of the gain tracking how low the
per-cell ceiling was rather than how hot the city is. A readiness
assessment should therefore report the scale at which land cover becomes
informative rather than a single per-city number.

Fourth, a diffusion--fidelity trade-off roughly twice as steep as a naive
estimate. Fifth, an evapotranspiration ablation showing that a standard
linear ET term contributes nothing to spatial-pattern skill at 500~m,
regardless of formulation.

The simulator itself achieves $r = +0.754$ on held-out MODIS days using
only inputs retrievable by anonymous HTTP request, which makes every
result above independently checkable without an institutional data
agreement.

%=================================================================
\vspace{6pt}

\funding{This research received no external funding.}

\institutionalreview{Not applicable.}

\informedconsent{Not applicable.}

\dataavailability{All input data are publicly available without
authentication: MODIS MYD11A1 and MOD13A1 via the Microsoft Planetary
Computer STAC API, ESA WorldCover v200 via the Copernicus/ESA object
store, ERA5 via the Google ARCO-ERA5 Zarr mirror, and administrative
boundaries via OpenStreetMap. Source code, calibration artifacts and the
scripts that regenerate every figure and table reported here are
available at \url{https://github.com/abhi-is-cool/austin-twin}.}

% REVIEW AND EDIT THIS TO MATCH WHAT YOU ACTUALLY DID -- MDPI requires
% disclosure of GenAI used for code, analysis, figures or text.
\acknowledgments{MODIS data are provided via the Microsoft Planetary
Computer; ESA WorldCover via Copernicus/ESA; ERA5 via Google ARCO-ERA5 on
Google Cloud Storage; OpenStreetMap data are copyright OpenStreetMap
contributors and licensed under the ODbL. During the preparation of this
manuscript, the author used Claude (Anthropic) for the purposes of
software implementation, analysis scripting, figure generation and
manuscript drafting. The author has reviewed and edited the output and
takes full responsibility for the content of this publication.}

\conflictsofinterest{The author declares no conflicts of interest.}

\abbreviations{Abbreviations}{%
The following abbreviations are used in this manuscript:\\

\noindent
\begin{tabular}{@{}ll}
ARCO & Analysis-Ready, Cloud-Optimised\\
CFL & Courant--Friedrichs--Lewy\\
ET & Evapotranspiration\\
LAI & Leaf Area Index\\
LST & Land Surface Temperature\\
NDVI & Normalised Difference Vegetation Index\\
PM & Penman--Monteith\\
STAC & SpatioTemporal Asset Catalog\\
UHI & Urban Heat Island\\
UTM & Universal Transverse Mercator\\
VPD & Vapour Pressure Deficit
\end{tabular}
}

%=================================================================
\begin{adjustwidth}{-\extralength}{0cm}

\reftitle{References}

\isAPAandChicago{}{%
\begin{thebibliography}{999}

\bibitem[Bhamjee et al.(2024)]{bhamjee2024detection}
Bhamjee, M.; et al. Detection and characterization of urban heat islands
with machine learning. In Proceedings of IGARSS 2024, pp. 1693--1699.

\bibitem[Bueno et al.(2013)]{bueno2013urban}
Bueno, B.; Norford, L.; Hidalgo, J.; Pigeon, G. The urban weather
generator. \textit{J. Build. Perform. Simul.} \textbf{2013}, \textit{6},
269--281.

% VERIFY: standard gradient-boosting reference, added for the nonlinear
% ceiling estimator. Check the volume and page range before submitting.
\bibitem[Friedman(2001)]{friedman2001greedy}
Friedman, J.H. Greedy function approximation: A gradient boosting machine.
\textit{Ann. Stat.} \textbf{2001}, \textit{29}, 1189--1232.

\bibitem[Hersbach et al.(2020)]{hersbach2020era5}
Hersbach, H.; et al. The ERA5 global reanalysis. \textit{Q. J. R.
Meteorol. Soc.} \textbf{2020}, \textit{146}, 1999--2049.

\bibitem[Hoang and Nguyen(2025)]{hoang2025geospatial}
Hoang, N.D.; Nguyen, Q.L. Geospatial analysis and machine learning
framework for urban heat island intensity prediction.
\textit{Sustainability} \textbf{2025}, \textit{17}, 4287.

\bibitem[HeatMat(2025)]{heatmat2025}
[Authors]. HeatMat: Simulation of city material impact on urban heat
island effect. \textit{arXiv} \textbf{2025}, arXiv:2601.22796.

\bibitem[Krayenhoff et al.(2021)]{krayenhoff2021albedo}
Krayenhoff, E.S.; et al. Cooling hot cities: A systematic and critical
review. \textit{Environ. Res. Lett.} \textbf{2021}, \textit{16}, 053007.

\bibitem[Kusaka et al.(2001)]{kusaka2001simple}
Kusaka, H.; et al. A simple single-layer urban canopy model.
\textit{Bound.-Layer Meteorol.} \textbf{2001}, \textit{101}, 329--358.

\bibitem[Lyu et al.(2022)]{lyu2022integrated}
Lyu, F.; et al. An integrated CyberGIS and machine learning framework for
fine-scale urban heat island prediction. \textit{Urban Inform.}
\textbf{2022}, \textit{1}, 6.

\bibitem[Masson(2000)]{masson2000physically}
Masson, V. A physically-based scheme for the urban energy budget.
\textit{Bound.-Layer Meteorol.} \textbf{2000}, \textit{94}, 357--397.

\bibitem[Microsoft(2023)]{microsoft_pc}
Microsoft. Planetary Computer STAC API. Available online:
\url{https://planetarycomputer.microsoft.com} (accessed on 1 September
2026).

\bibitem[Mitraka et al.(2025)]{mitraka2025prototypes}
Mitraka, Z.; et al. Prototypes for enhanced urban heat monitoring.
\textit{Discov. Cities} \textbf{2025}, \textit{2}, 14.

\bibitem[Oke(1982)]{oke1982energetic}
Oke, T.R. The energetic basis of the urban heat island. \textit{Q. J. R.
Meteorol. Soc.} \textbf{1982}, \textit{108}, 1--24.

\bibitem[Santamouris(2014)]{santamouris2014cool}
Santamouris, M. Cooling the cities---A review of reflective and green
roof mitigation technologies. \textit{Sol. Energy} \textbf{2014},
\textit{103}, 682--703.

\bibitem[Storn and Price(1997)]{storn1997differential}
Storn, R.; Price, K. Differential evolution---A simple and efficient
heuristic. \textit{J. Glob. Optim.} \textbf{1997}, \textit{11}, 341--359.

\bibitem[Voogt and Oke(2003)]{voogt2003thermal}
Voogt, J.A.; Oke, T.R. Thermal remote sensing of urban climates.
\textit{Remote Sens. Environ.} \textbf{2003}, \textit{86}, 370--384.

\bibitem[Wang and Li(2022)]{wang2022berlin}
Wang, X.; Li, H. The effectiveness of cool and green roofs in mitigating
urban heat island. \textit{Build. Environ.} \textbf{2022}, \textit{217},
109082.

\bibitem[Wesley et al.(2026)]{wri2026cooling}
Wesley, E.J.; Mackres, E.; Pool, J. Cooling potential of urban trees.
\textit{WRI Insights} \textbf{2026}.

\bibitem[Zanaga et al.(2022)]{zanaga2022esa}
Zanaga, D.; et al. ESA WorldCover 10~m 2021 v200. \textit{Zenodo}
\textbf{2022}.

\bibitem[Zhou et al.(2025)]{zhou2025tree}
Zhou, W.; et al. Increasing tree canopy lowers urban air temperature by
up to 1.5~$^{\circ}$C. \textit{npj Urban Sustain.} \textbf{2025},
\textit{5}, 92.

\bibitem[Ziter et al.(2019)]{ziter2019scale}
Ziter, C.D.; et al. Scale-dependent interactions between tree canopy
cover and impervious surfaces. \textit{Proc. Natl. Acad. Sci. USA}
\textbf{2019}, \textit{116}, 7575--7580.

\end{thebibliography}
}

\PublishersNote{}
\end{adjustwidth}
\end{document}
